\section{Algorithmic Details and Theorem-to-Code Map}
\label{app:algorithms}

\subsection{Tokenizer pseudocode}

The tokenizer implemented in \codepath{mino/models/wavepacket.py} follows the steps below.

\begin{enumerate}
\item unfold the input field into overlapping patches;
\item apply a window function to each patch;
\item compute a real local \FFT\ on each patch;
\item retain a fixed set of local modes ordered by increasing radius;
\item package the retained coefficients together with position, frequency, and scale metadata.
\end{enumerate}

This design matters because the theorem is written over finite packet coefficients. The tokenizer is the bridge from grid data to the theorem's coefficient space.

\subsection{Propagation-layer pseudocode}

The propagation block implemented in \codepath{mino/models/layers.py} can be summarized as:
\begin{enumerate}
\item concatenate packet features with metadata;
\item use an MLP to predict a small update of the packet metadata;
\item build a metadata-space distance matrix after the update;
\item mask that kernel to the retained top-$K$ metadata neighbors;
\item aggregate local messages through the masked kernel induced by those distances;
\item gate the aggregated message and add it residually to the packet features.
\end{enumerate}
The backward-compatible implementation computes the pairwise distance matrix densely, but the actual interaction weights are top-$K$ truncated before aggregation. The theorem-aligned profile also supports sparse aggregation after the top-$K$ set has been selected. Thus the model matches the finite bounded-stencil hypothesis at the level of retained interactions; true sparse-neighbor search is the speed and memory optimization.

\subsection{Symbol branch}

The generic empirical symbol branch is another MLP that reads packet features and updated metadata and outputs a scale and shift. The shift is interpreted as a local empirical source or residual correction, not as the theorem-facing symbol law. In the theorem-aligned \texttt{spectral\_order} mode, the branch is multiplier-only: it reads phase-space metadata, log radial frequency, and normalized frequency direction, applies the configured order weight, and multiplies packet coefficients without an additive shift. This is the finite \(S^m\)-proxy used by the mechanism profile. Its main function is to separate amplitude correction from transport while keeping the theorem-facing branch local and packet-diagonal.

\subsection{Low-frequency residual branch}

The low-frequency branch is a compact spectral residual block with zero initialization. It is intentionally conservative at initialization, so that the model starts close to an identity skip path rather than injecting arbitrary spectral noise. This design was chosen precisely because the tiny pilot benchmark exposed instability when the spectral residual and skip branch were both randomly initialized.

\subsection{Why the initialization matters}

The initial version of the code used a random skip map and a random low-frequency spectral block. This created extremely poor one-epoch pilot behavior for \MiNO. The current version initializes the spectral residual to zero and the skip path to an identity-like map, which is the correct theorem-aware choice:
\begin{itemize}
\item the residual branch should start neutral,
\item the transported-symbol branch should learn departures from a stable baseline,
\item the finite theorem budgets should not be polluted by arbitrary initialization noise.
\end{itemize}

\subsection{Theorem-to-code map}

The correspondence between the manuscript and the repository is as follows.

\begin{center}
\begin{tabular}{p{0.30\textwidth}p{0.58\textwidth}}
\toprule
Mathematical object & Implementation site \\
\midrule
finite packet coefficients & \codepath{WavepacketEncoding.features} \\
packet metadata $(x,\xi,\sigma)$ & \codepath{WavepacketEncoding.metadata} \\
reference / approximate transport index & implicit in learned metadata update + kernel aggregation \\
approximate symbol & \codepath{SymbolModulationLayer} \\
residual branch & low-frequency spectral residual + synthesis mismatch \\
packet-budget diagnostics & \codepath{mino/metrics/packet\_bridge.py} \\
Navier--Stokes rollout loaders & \codepath{mino/data/rollout.py} \\
autoregressive rollout evaluator & \codepath{mino/training/rollout.py} \\
pointwise propagation theorem & Lean theorem \codepath{pointwise\_propagation\_bound} \\
$\ell^1$ propagation theorem & Lean theorem \codepath{l1\_propagation\_bound} \\
approximation corollary & Lean theorem \codepath{finite\_approximation\_corollary} \\
\bottomrule
\end{tabular}
\end{center}

\subsection{Implementation scope}

The theory-driven decomposition is present, with several intentionally simple implementation choices:
\begin{enumerate}
\item the local packet frame is local Fourier/Gaussian instead of a richer curvelet, shearlet, or adaptive semiclassical system;
\item top-$K$ interactions are selected after dense distance evaluation, with sparse aggregation available after selection;
\item the benchmark data include synthetic surrogates and cached PDE families;
\item the proxy objective is a diagnostic packet matching loss tied to theorem budgets.
\end{enumerate}
These choices define the present empirical scope.

\subsection{Practical next steps}

The next technical milestones are clear from the theorem-to-code map:
\begin{enumerate}
\item replace dense pairwise distance evaluation by sparse or approximate-neighbor top-$K$ transport selection;
\item enrich the packet basis beyond local Fourier patches;
\item add theorem-aware auxiliary losses for transport and symbol supervision;
\item connect the benchmark harness to full PDE datasets and official baselines.
\end{enumerate}
The manuscript therefore provides the theorem-code foundation and the engineering roadmap in one artifact.

\subsection{End-to-end forward-pass pseudocode}

For completeness, we record the entire forward pass of the current \MiNO\ scaffold in theorem-aware order.

\begin{enumerate}
\item analyze the input field into packet coefficients and metadata;
\item lift packet coefficients to hidden features;
\item repeat for each microlocal block:
\begin{enumerate}
\item predict a metadata update from current packet features,
\item apply local aggregation in the updated metadata geometry,
\item modulate the aggregated packet features with the symbol branch,
\item add the result to the running packet features;
\end{enumerate}
\item decode the packet features back to coefficient space;
\item synthesize a field from the decoded coefficients;
\item add the low-frequency residual branch in physical space.
\end{enumerate}

This order is important. If one interleaves the low-frequency branch too early, the model can short-circuit the packet transport path and the theorem-aware decomposition becomes harder to interpret.

\subsection{Estimating the theorem budgets in practice}

The theorem is stated with budgets $(\eps_{\mathrm{tr}},\eps_{\mathrm{sym}},\eps_{\mathrm{res}})$. In code, a practical estimator for those quantities can be organized as follows.

\paragraph{Transport estimate.}
Compute packet clouds for the source and target fields. Restrict to the dominant packet mass and match source-to-target packets by nearest metadata distance or by an optimal-transport assignment. The average matched distance is the transport proxy.

\paragraph{Symbol estimate.}
For matched packets, compute the ratio between target and transported source amplitudes. Compare that ratio to the symbol branch output.

\paragraph{Residual estimate.}
Apply the current learned transport and symbol to the source packet cloud and measure unmatched packet energy in the target cloud.

\paragraph{Envelope estimates.}
The constants $B_a$ and $B_s$ should be reported not as arbitrary global constants, but as observed maxima or quantiles over the packet statistics used in a given benchmark family.

This procedure supplies theorem-aware diagnostics beyond the final field norm.

\subsection{Sparse transport stencil design}

The propagation block is logically stencil-truncated and uses a dense neighbor search: it computes all pairwise metadata distances and then selects the retained top-$K$ neighbors. The \texttt{sparse\_topk} path makes the aggregation after that selection sparse. The computational bottleneck is the dense search used to find the stencil; a sparse transport graph removes that bottleneck.

\paragraph{Stencil proposal.}
For each packet, retain only the $k$ nearest metadata neighbors after the predicted transport update, or retain all neighbors inside a fixed metadata radius. Then perform message aggregation only on that restricted graph.

\paragraph{Why this matches the theory.}
The continuous reduction theorem says that packet interactions decay away from the canonical graph. A sparse transport stencil is therefore not merely a computational optimization; it is the discrete numerical image of the off-graph decay assumption.

\paragraph{Expected effect.}
This change should:
\begin{itemize}
\item reduce quadratic cost,
\item reduce noisy long-range packet interactions,
\item make the transport branch more interpretable,
\item improve alignment with the continuous packet matrix picture.
\end{itemize}

\subsection{Training and rollout loop pseudocode}

A theorem-aware training loop for the flagship version should look like:
\begin{enumerate}
\item sample a minibatch of operator pairs together with packet-derived auxiliary diagnostics;
\item run the \MiNO\ forward pass;
\item compute field loss, transport proxy loss, symbol proxy loss, and residual loss;
\item apply branch-specific weighting or curriculum scheduling;
\item update the model;
\item periodically evaluate ablation toggles and OOD diagnostics.
\end{enumerate}
The current code instantiates this loop for field-only and proxy-augmented one-step training. A separate rollout evaluator then repeatedly composes the learned one-step map, records field metrics, and records packet-budget diagnostics at each step. This separation is important: one-step training remains the controlled approximation problem, while rollout is the nonlinear accumulation stress test.

\subsection{Theorem-to-diagnostic map}

It is useful to record explicitly which diagnostics correspond to which theorem terms.

\begin{center}
\begin{tabular}{p{0.32\textwidth}p{0.28\textwidth}p{0.26\textwidth}}
\toprule
Theorem term & Practical estimator & Most diagnostic setting \\
\midrule
$\eps_{\mathrm{tr}}$ & packet matching or metadata OT cost & wave, Helmholtz \\
$\eps_{\mathrm{sym}}$ & post-match amplitude ratio error & all regimes \\
$\eps_{\mathrm{res}}$ & unmatched packet energy & mixed and oscillatory regimes \\
$B_a$ & packet coefficient envelope statistic & all regimes \\
$B_s$ & observed symbol envelope statistic & all regimes \\
\bottomrule
\end{tabular}
\end{center}

This map keeps the theorem and the experiments in the same diagnostic language.

\subsection{From local Fourier patches to richer packet systems}

The tokenizer is intentionally simple. The representational upgrade path is:

\paragraph{Wavepacket upgrade.}
Replace pure patch-local Fourier modes by atoms with explicit scale-frequency anisotropy and smoother spatial tapering.

\paragraph{Curvelet upgrade.}
Adopt a representation better adapted to wavefront geometry and anisotropic singularities. This is the mathematically most natural direction for wave propagation, though it raises implementation complexity.

\paragraph{Semiclassical/FBI-style upgrade.}
Introduce a packetization closer to the continuous microlocal transform used in the reduction theorem. This would tighten the connection between code and theory the most.

Each upgrade reduces the gap between the executable packetizer and the continuous packet matrix picture of Appendix~\ref{app:continuous}.

\subsection{Implementation invariants}

Regardless of the eventual tokenizer or transport stencil, several invariants should remain fixed.
\begin{enumerate}
\item packet analysis should remain mathematically interpretable;
\item transport and symbol should remain distinct branches;
\item the low-frequency residual should remain a secondary path rather than the main explanatory path;
\item ablations should always be able to switch these components on and off cleanly.
\end{enumerate}
These invariants are more important than any one local coding decision. They are what preserve the single-spine theorem structure of the paper even as the implementation grows more sophisticated.
